\documentclass{article}
\usepackage{graphicx} % Required for inserting images
\usepackage[margin=2cm, top=3cm]{geometry} % margine superiore ridotto a 1 cm
\usepackage{listings}
\usepackage{subcaption}  % Pacchetto per sottodidascalie
\usepackage{float}       % Pacchetto per posizionamento rigido delle figure
\usepackage{xcolor} % Pacchetto per i colori
\usepackage{amsmath} % Pacchetto per simboli avanzati

\title{Domande Economia Applicata all'Ingegneria}
\author{Mattia Cosimi}
\date{December 2024}

\begin{document}

\maketitle

\section{Prima parte}
\subsubsection{Si descriva formalmente il problema del costo minimo di produzione}
\subsubsection{Funzioni di produzione caratterizzate da rendimenti di scala costanti possono condurre a costi caratterizzati da economie di scala? }
\subsubsection{Che cosa si intende per curva di indifferenza nella teoria delle preferenze?}
\subsubsection{Descrivere formalmente il problema della scelta ottima di un consumatore}
\subsubsection{Si illustrino gli effetti sul vincolo di bilancio di politiche di tassazione sul reddito, sui prezzi e di razionamento}
\subsubsection{Si illustri la condizione di ottimalità interna del problema della scelta ottima del consumatore}
\subsubsection{Si definiscano i concetti di funzione di produzione, di isoquanti di produzione e di tasso marginale di sostituzione tecnica}
\subsubsection{Quando due beni sono sostituti e quando sono complementi? Cosa si intende per perfette sostitutibilità e complementarità}
\subsubsection{Quando un paniere è preferito ad un altro e quando è indifferente? Cosa rappresenta la curva di indifferenza dei panieri e come si determina a partire dalla funzione di utilità? Fornire un esempio illustrativo}
\subsubsection{Illustrare la determinazione della curva di costo di lungo periodo a partire dalla soluzione del problema di costo minimo}
\subsubsection{Illustrare il generico problema del costo minimo, dati w1,w2, f(x1,x2), y. Come cambia il problema se siamo nel breve periodo ed un fattore di produzione è stato già acquistato sotenendo, quindi, un costo fisso pari a F?}

\end{document}
